% ============================================================
\chapter*{Preface}
\addcontentsline{toc}{chapter}{Preface}
% ============================================================

Machine learning is a discipline at the intersection of mathematics,
computer science, and statistics. Its fundamental objective is to design
algorithms capable of \textbf{learning from data}: instead of explicitly
programming a decision rule, we let the algorithm discover this rule from
examples.

\section*{Course Positioning}

This course is at the L3/Master~1 level and assumes prerequisites in:
\begin{itemize}
  \item \textbf{Linear algebra}: vector spaces, matrices, eigenvalues,
    singular value decomposition.
  \item \textbf{Probability and statistics}: random variables, expectation,
    variance, normal distribution, estimation, tests.
  \item \textbf{Analysis}: partial derivatives, gradient, optimization of
    multivariate functions.
  \item \textbf{Python programming}: NumPy, Matplotlib; knowledge of
    scikit-learn is a plus but not required.
\end{itemize}

\section*{Pedagogical Philosophy}

Each chapter follows a three-step progression:
\begin{enumerate}
  \item \textbf{Theory}: rigorous mathematical derivation of algorithms,
    with statement and proof of key results.
  \item \textbf{Practice}: implementation in Python with scikit-learn and,
    when instructive, from scratch.
  \item \textbf{Analysis}: discussion of strengths, weaknesses, assumptions,
    and use cases of each method.
\end{enumerate}

\section*{Course Outline}

\begin{center}
\begin{tikzpicture}[
  block/.style={rectangle, draw, rounded corners, minimum width=4cm,
    minimum height=0.8cm, align=center, font=\small},
  arr/.style={-{Stealth[length=2.5mm]}, thick, steelblue}
]
  \node[block, fill=blue!10] (ch1) {Ch.~1: Introduction};
  \node[block, fill=blue!10, below=0.3cm of ch1] (ch2) {Ch.~2: Regression};
  \node[block, fill=blue!10, below=0.3cm of ch2] (ch3) {Ch.~3: Classification};
  \node[block, fill=green!10, below=0.3cm of ch3] (ch4) {Ch.~4: Regularization};
  \node[block, fill=green!10, below=0.3cm of ch4] (ch5) {Ch.~5: SVM};
  \node[block, fill=green!10, below=0.3cm of ch5] (ch6) {Ch.~6: Trees};

  \node[block, fill=orange!10, right=2cm of ch1] (ch7) {Ch.~7: Ensembles};
  \node[block, fill=orange!10, below=0.3cm of ch7] (ch8) {Ch.~8: Unsupervised};
  \node[block, fill=orange!10, below=0.3cm of ch8] (ch9) {Ch.~9: Dim.\ Reduction};
  \node[block, fill=red!10, below=0.3cm of ch9] (ch10) {Ch.~10: Bayesian};
  \node[block, fill=red!10, below=0.3cm of ch10] (ch11) {Ch.~11: Selection};
  \node[block, fill=red!10, below=0.3cm of ch11] (ch12) {Ch.~12: Kernels};

  \draw[arr] (ch1) -- (ch2);
  \draw[arr] (ch2) -- (ch3);
  \draw[arr] (ch3) -- (ch4);
  \draw[arr] (ch4) -- (ch5);
  \draw[arr] (ch5) -- (ch6);
  \draw[arr] (ch6.east) -- ([xshift=0.5cm]ch6.east) |- (ch7.west);
  \draw[arr] (ch7) -- (ch8);
  \draw[arr] (ch8) -- (ch9);
  \draw[arr] (ch9) -- (ch10);
  \draw[arr] (ch10) -- (ch11);
  \draw[arr] (ch11) -- (ch12);
\end{tikzpicture}
\end{center}

\noindent Chapters 1--6 cover the foundations (linear and nonlinear
supervised learning). Chapters 7--9 address ensemble methods and
unsupervised learning. Chapters 10--12 treat advanced topics (Bayesian
approach, model selection, kernel methods).

\vfill
\noindent\textit{Happy reading and happy learning!}
